%! Comparing Studies and Choosing a Method — Unit 1, Lesson 1.7 — Guided Notes
\documentclass[10pt]{article}
\usepackage{saar-article}
\usepackage{saar-boxes}
\newcommand{\TallMath}[1]{$\displaystyle #1\rule[-1.4em]{0pt}{3.2em}$}

\begin{document}

\pageheader{Unit 1, Lesson 1.7}{Guided Notes}
% NAMESTRIP: no \namedateperiod here. The student writes their name once, on the
% cover this component is stapled behind. See references/conventions.md.

\vspace{0.15in}

\begin{objectivebox}
\small
By the end of this lesson, I will be able to\dots
\begin{enumerate}[leftmargin=1.5em, itemsep=2pt, topsep=2pt]
  \item compare an \textbf{observational study} with a \textbf{controlled
        experiment}, point by point;
  \item state the strongest conclusion each one earns, and \emph{justify} the
        difference;
  \item decide when an experiment is \textbf{impossible or unethical} to run,
        and say what to run instead;
  \item use the \textbf{statistical cycle} to plan a well-designed experiment;
  \item choose a \textbf{data collection method} --- survey, interview, focus
        group, observation, or content analysis --- that fits a given context.
\end{enumerate}
\end{objectivebox}

\boxguard
\begin{vocabbox}
\small
Fill in each term as we build it together.
\termblanklong{Controlled experiment}
\termblanklong{Observational study}
\termblanklong{Survey}
\termblanklong{Interview}
\termblanklong{Focus group}
\termblanklong{Content analysis}
\end{vocabbox}

\boxguard[12]
\begin{hookbox}
\small
Riverbend's counselor now has both studies in front of her. Watching $150$
records said $65\%$ against $80\%$ --- the Study Center looked harmful.
Assigning $120$ volunteers said $75\%$ against $55\%$ --- it looked helpful. You
already know which one she should believe.

\vspace{2pt}
Now the principal brings a different question: \emph{do students who ride the
bus more than $45$ minutes each way score lower on the state test?}

\vspace{2pt}
\textbf{Could the counselor settle this one with an experiment the way she
settled the Study Center? What exactly would she have to assign?}
\writelines{2}
\end{hookbox}

\boxguard[30]
\begin{notesbox}{1. The Two Study Types, Side by Side}
\small
Both study types compare groups. They differ in exactly one place --- and that
one place decides everything you are allowed to say at the end.

\begin{center}
\renewcommand{\arraystretch}{1.4}
\begin{tabularx}{0.98\linewidth}{@{} p{3.1cm} p{4.5cm} X @{}}
\toprule
\textbf{Ask this} & \textbf{Observational study}
  & \textbf{Controlled experiment} \\
\midrule
Who decides which group each individual is in?
  & The individuals themselves, or something that already happened to them.
  & \blank{3.8cm} \\
Does the researcher impose a treatment?
  & No --- the researcher only measures or asks.
  & \blank{3.8cm} \\
Are the two groups alike before the study starts?
  & \blank{3.8cm} & \blank{3.8cm} \\
Strongest conclusion it earns
  & \blank{3.8cm} & \blank{3.8cm} \\
\bottomrule
\end{tabularx}
\end{center}

\vspace{3pt}
Now put Riverbend's two studies in that frame.

\begin{center}
\renewcommand{\arraystretch}{1.45}
\begin{tabularx}{0.98\linewidth}{@{} p{4.1cm} p{3.6cm} X @{}}
\toprule
\textbf{The study} & \textbf{What it found}
  & \textbf{The strongest sentence it earns} \\
\midrule
Watched $150$ records (1.5) & $65\%$ of users vs.\ $80\%$ of non-users
  & \blank{4.3cm} \\
Assigned $120$ volunteers (1.6) & $75\%$ treated vs.\ $55\%$ control
  & \blank{4.3cm} \\
\bottomrule
\end{tabularx}
\end{center}

\vspace{2pt}
\textit{An experiment is the stronger design. The next section is about why you
cannot always have one.}
\end{notesbox}

\boxguard[30]
\begin{notesbox}{2. When You Cannot Run an Experiment}
\small
To run an experiment you must be able to \textbf{assign} the treatment. Three
things can stop you:

\begin{enumerate}[leftmargin=1.6em, itemsep=2pt, topsep=3pt]
  \item \textbf{Unethical} --- assigning the treatment would harm someone.
  \item \textbf{Impossible} --- the ``treatment'' is something a person already
        \emph{is} or \emph{has}. You cannot hand it out.
  \item \textbf{Impractical} --- it would cost too much, take too long, or
        nobody would agree to be assigned.
\end{enumerate}

\vspace{3pt}
For each question, decide whether the researcher could assign the groups.

\begin{center}
\renewcommand{\arraystretch}{1.4}
\begin{tabularx}{0.98\linewidth}{@{} X c p{4.5cm} @{}}
\toprule
\textbf{The question} & \textbf{Can you assign?}
  & \textbf{If not, why --- and what do you run instead?} \\
\midrule
Does a $10$-minute silent reading block raise reading scores?
  & \blank{1.1cm} & \blank{4.1cm} \\
Do students who smoke have worse lung function at $30$?
  & \blank{1.1cm} & \blank{4.1cm} \\
Do retired pool members sleep more than members who work full time?
  & \blank{1.1cm} & \blank{4.1cm} \\
Do students who ride the bus over $45$ minutes score lower?
  & \blank{1.1cm} & \blank{4.1cm} \\
\bottomrule
\end{tabularx}
\end{center}

\vspace{3pt}
When you cannot assign, you run the best observational study you can --- and you
change the \emph{verb}, not the data. A study like that may report
a(n) \blank{3.4cm}, but never a \blank{3.4cm}.

\vspace{2pt}
\textit{``We could not run an experiment'' is not an excuse for a stronger
claim. It is the reason for a weaker one.}
\end{notesbox}

\boxguard[22]
\begin{notesbox}{3. Planning an Experiment with the Statistical Cycle}
\small
Lesson 1.5 walked the four stages for an observational study. Here is the same
cycle for an \emph{experiment}. The counselor's new question: does a weekly
one-on-one check-in raise the share of students who turn in every assignment?

\begin{center}
\renewcommand{\arraystretch}{1.45}
\begin{tabularx}{0.98\linewidth}{@{} p{2.6cm} p{4.6cm} X @{}}
\toprule
\textbf{Stage} & \textbf{What the stage asks} & \textbf{Our check-in experiment} \\
\midrule
Formulate & What do we want to know, and about whom?
  & Does a weekly check-in raise the share of students who turn in every
    assignment? \\
Collect & Who are the units, and how does \emph{chance} put them in groups?
  & \blank{4.4cm} \\
Represent & How will we organize what comes back?
  & \blank{4.4cm} \\
Analyze \& report & What do the numbers say, \emph{in context}?
  & \blank{4.4cm} \\
\bottomrule
\end{tabularx}
\end{center}

\vspace{3pt}
One stage does something Lesson 1.5's plan never did. Which stage is it, and
what does it do? \blank{9cm}

\vspace{3pt}
That stage is also where the four principles get checked: two groups
(\textbf{comparison}), a generator (\textbf{randomization}), plenty of students
in each group (\textbf{replication}), and everything else held the same
(\textbf{control}).

\vspace{2pt}
\textit{Same four stages every time. What an experiment adds is one word inside
stage two: \emph{assign}.}
\end{notesbox}

\boxguard[30]
\begin{notesbox}{4. Five Ways to Collect the Data}
\small
Once you know which study you are running, you still have to decide \emph{how}
the data arrive. These five are the standard menu.

\begin{center}
\renewcommand{\arraystretch}{1.35}
\begin{tabularx}{0.98\linewidth}{@{} p{2.7cm} X @{}}
\toprule
\textbf{Method} & \textbf{Best when you want\dots} \\
\midrule
Survey & the same fixed questions answered by \emph{many} people --- counts,
  percents, attitudes across a whole sample. \\
Interview & \emph{depth} from one person at a time --- the reasons behind an
  answer, in that person's own words. \\
Focus group & a small group talking \emph{together}, where one person's answer
  reminds the next person of theirs. \\
Observation & to record what people actually \emph{do}, when what they do may
  not match what they say. \\
Content analysis & to study records or documents that \emph{already exist} ---
  useful for change over time. \\
\bottomrule
\end{tabularx}
\end{center}

\vspace{3pt}
Match each context to the method that fits it best.

\begin{center}
\renewcommand{\arraystretch}{1.4}
\begin{tabularx}{0.98\linewidth}{@{} X p{4cm} @{}}
\toprule
\textbf{What the researcher needs} & \textbf{Best method} \\
\midrule
The library wants five years of borrowing records summarized.
  & \blank{3.6cm} \\
The pool wants to know how many of its $1{,}500$ members would use a
$6$~a.m.\ lane. & \blank{3.6cm} \\
The counselor wants to understand \emph{why} four students stopped coming to
the Study Center. & \blank{3.6cm} \\
The cafeteria wants to know whether students actually take the fruit --- not
whether they say they do. & \blank{3.6cm} \\
The library wants eight teens to react to three new logos and build on each
other's comments. & \blank{3.6cm} \\
\bottomrule
\end{tabularx}
\end{center}

\vspace{3pt}
\textbf{Careful:} \emph{observation} is a way to \emph{collect} data;
an \emph{observational study} is a kind of \emph{study}. In Bayside's reading
experiment a teacher watched and recorded how many minutes each student actually
read, so the method was \blank{3.2cm} --- and the study was still
a(n) \blank{3.6cm}, because chance assigned the groups.

\vspace{2pt}
\textit{The study type is decided by who built the groups. The method is decided
by what the question needs.}
\end{notesbox}

\boxguard
\begin{practicebox}
\small
\textbf{Cedar Ridge Apartments.} The manager of a $600$-household complex has
three questions this month. For each one, decide which study she can run and how
she should collect the data.

\begin{center}
\renewcommand{\arraystretch}{1.45}
\begin{tabularx}{0.98\linewidth}{@{} X p{2.6cm} p{3.1cm} @{}}
\toprule
\textbf{Her question} & \textbf{Experiment or observational?}
  & \textbf{Best collection method} \\
\midrule
Does a weekly recycling reminder increase how much a household recycles?
  & \blank{2.2cm} & \blank{2.7cm} \\
Do households in the renovated building recycle more than households in the old
building? & \blank{2.2cm} & \blank{2.7cm} \\
Why do some households never use the recycling bins at all?
  & \blank{2.2cm} & \blank{2.7cm} \\
\bottomrule
\end{tabularx}
\end{center}

\vspace{3pt}
She finds that the renovated building does recycle more. Write the strongest
sentence she may publish about that result, and name one other variable that
could explain the difference instead.
\writelines{2}
\end{practicebox}

\end{document}
